\naslov{Relativno gibanje}

\begin{definicija}
  Koordinatni sistem $\varphi(t, \mb{P})$ se \pojem{giblje} glede na koordinatni
  sistem $\varphi'(t', \mb{P}')$, če obstaja trojica $(\mb{P}_0, \mb{P}_0', Q)$,
  kjer je $\mb{P}_0 \in \mathbb{E}$, $\mb{P}_0' : \R \to \mathbb{E}$, in $Q: \R
  \to SO(3)$, da velja $t = t'$ in
  \[
	\mb{P}'(t) = \mb{P}_0'(t) + Q(t) (\mb{P} - \mb{P}_0).
  \]
\end{definicija}

Predpostavimo, da je $\varphi'$ inercialen, in ga imenujmo \pojem{absolutni
  koordinatni sistem}, $\varphi$ pa je \pojem{relativni koordinatni sistem}.

\begin{trditev}
  Rotacijski del gibanja je neodvisen od izbire trojice.
\end{trditev}

\begin{proof}
  Recimo
  \[
	\mb{P}' = \mb{P}_0' + Q (\mb{P} - \mb{P}_0)
	= \tilde{\mb{P}}_0' + \tilde{Q} (\mb{P} - \tilde{\mb{P}}_0).
  \]
  Za $\mb{P} = \mb{P}_0$ dobimo
  \[
	\mb{P}_0' = \tilde{\mb{P}}_0' + \tilde{Q} (\mb{P}_0 - \tilde{\mb{P}}_0).
  \]
  Če to vstavimo nazaj gor, pridemo do
  \[
	Q(\mb{P} - \mb{P}_0) = \tilde{Q}(\mb{P} - \mb{P}_0),
  \]
  torej $Q = \tilde{Q}$.
\end{proof}

\vprasanje{Kdaj se koordinatni sistem giblje glede na nek drugi koordinatni
  sistem? Dokaži, da je rotacijski del neodvisen od izbire trojice.}

Za odvod velja
\[
  \pvec{v}' = \dot{\mb{P}}_0' + \dot{Q} (\mb{P} - \mb{P}_0) + Q \dot{\mb{P}}
  = Q(Q^T \pvec{v}_0' + Q^T \dot{Q} (\mb{P} - \mb{P}_0) + \vec{v}_{\text{rel}}).
\]
Prvemu členu pravimo \pojem{translatorna hitrost}, drugemu \pojem{rotacijska
  hitrost}, tretjemu pa \pojem{relativna hitrost}.

\begin{trditev}
  $Q^T \dot{Q}$ je poševno simetrični tenzor.
\end{trditev}

\begin{proof}
  Velja $Q^T Q = I$.
  Če to odvajamo, dobimo
  \[
	\partial_t (Q^T) Q + Q^T \dot{Q} = 0.
  \]
\end{proof}

\begin{izrek}
  Naj bo $W$ poševno simetričen na trirazsežnem evklidskem prostoru.
  Potem obstaja vektor $\vec{\omega}$ tako, da je $W \vec{a} = \vec{\omega}
  \times \vec{a}$ za vsak $\vec{a}$.
\end{izrek}

\begin{proof}
  Vemo, da obstaja lastna vrednost $W \vec{p} = \lambda \vec{p}$.
  Ker je
  \[
	\lambda \abs{\vec{p}}^2 = \vec{p} \cdot W \vec{p}
	= W^T \vec{p} \cdot \vec{p} = - \lambda \abs{\vec{p}}^2,
  \]
  torej $\lambda = 0$.

  BŠS je $\abs{\vec{p}} = 1$.
  Dopolnimo ga lahko do ortonormirane baze prostora $\{\vec{p}, \vec{q},
  \vec{r}\}$, kjer je $\vec{r} = \vec{p} \times \vec{q}$.
  Ker je $\vec{a} W \vec{a} = 0$, dobimo
  \begin{align*}
	W \vec{q} = W \vec{q} \cdot \vec{r} . \vec{r},
	&& W \vec{r} = W \vec{r} \cdot \vec{q} . \vec{q}.
  \end{align*}
  Za poljuben $\vec{a}$ je
  \[
	W \vec{a} = \vec{a} \cdot \vec{q} . W \vec{q} \cdot \vec{r} . \vec{r}
	+ \vec{a} \cdot \vec{r} . W \vec{r} \cdot \vec{q} . \vec{q}
  \]
  Velja $\vec{q} \times \vec{r} = \vec{p}$ in $\vec{r} \times \vec{p} =
  \vec{q}$, torej
  \[
	W \vec{a} = W \vec{q} \cdot \vec{r} . (\vec{a} \cdot \vec{q} . \vec{p}
	\times \vec{q} + \vec{a} \cdot \vec{r} . \vec{p} \times \vec{r})
	= W \vec{q} \cdot \vec{r} . \vec{p} \times \vec{a}.
  \]
\end{proof}

\vprasanje{Dokaži, da poševno simetričen tenzor deluje kot vektorski produkt.}
